%====================================================================
% Introdução: Escreva logo após o \chapter{Introdução} a texto 
% de seu trabalho referente a introdução.
%====================================================================
\chapter{Introdução}
O escoamento de filmes finos é um fenômeno quase onipresente no cotidiano comum e industrial. Podendo ser encontrado em fenômenos biofísicos, como a formação de lágrimas e a lubrificação de articulações, em processos industriais, como o transporte de petróleo pesado em tubulações e os ditos filmes de evaporação, e em fenômenos naturais, como a formação de gotas de chuva e a erosão de rochas. Dado a sua ubiquidade, este fenômeno vem sendo estudado há décadas, com o objetivo de compreender e manipular sua dinâmica \citep{Oron1997, Craster2009}.

Apesar de não se ter um consenso sobre a definição de filme fino, algumas áreas possuem um consenso na literatura, por exemplo, o escoamento de filmes finos de fluído, começou a ser mais estudado após o trabalho de \cite{reynolds1886} sobre lubrificação. A partir desse trabalho, uma das possíveis definições de filme fino para este caso específico é $\epsilon << 1$, no qual $\epsilon = h/L$ na qual $h$ é a espessura do filme e $L$ é o comprimento característico do escoamento.

Tratando-se de escoamentos onde o filme fino é composto por um líquido, essa definição demonstra uma das maiores dificuldades no estudo desse fenômeno quanto à simulação computacional do mesmo: a disparidade entre as grandezas características entre o filme e o escoamento faz com que a malha computacional necessária seja, muitas vezes, proibitiva. Isto faz com que métodos clássicos de captura direta de interface, como o método VOF, não sejam viáveis, devido a necessidade de uma malha extremamente refinada em relação ao domínio como um todo \citep{Kakimpa2016, Kakimpa2015, Morris2017}. Apesar disto, vários trabalhos vêm sendo desenvolvidos para tentar contornar essa dificuldade, como os métodos reduzidos, que buscam modelar o filme fino de forma simplificada, e os métodos híbridos, que buscam acoplar um modelo reduzido com um modelo de captura direta de interface.

